\documentclass[11pt,a4paper]{article}
\usepackage[utf8]{inputenc}
\usepackage[T1]{fontenc}
\usepackage[catalan]{babel}
\usepackage{lmodern}
\usepackage[sfdefault,lining]{FiraSans}
\usepackage[a4paper,top=2.3cm,bottom=2.2cm,left=1.8cm,right=1.8cm,headheight=1.3cm,headsep=0.4cm,footskip=1cm]{geometry}
\usepackage{amsmath,amssymb}
\usepackage{icomma}
\usepackage{graphicx}
\usepackage[table]{xcolor}
\usepackage{tikz}
\usepackage{fontawesome5}
\usepackage{ccicons}
\usepackage{fancyhdr}
\usepackage{lastpage}
\usepackage{tasks}
\usepackage{needspace}
\usepackage[most]{tcolorbox}
\usepackage{array}
\usepackage[hidelinks,pdftitle={Solucions de la fitxa de la unitat 6: Sistemes d'equacions. 3r d'ESO},pdfsubject={Matemàtiques 3r ESO},pdfkeywords={sistemes d'equacions, equacions lineals amb dues incògnites, substitució, igualació, reducció, mètode gràfic, classificació de sistemes, problemes}]{hyperref}

\definecolor{navy}{HTML}{1B2748}
\definecolor{accent}{HTML}{E8603A}
\definecolor{teal}{HTML}{1F9E90}
\definecolor{soft}{HTML}{EEF1F8}
\definecolor{muted}{HTML}{5D6685}

\setlength{\parindent}{0pt}
\setlength{\parskip}{3pt}
\settasks{label=\alph*), label-format=\bfseries\color{navy}, label-width=1.3em, item-indent=1.8em, column-sep=0.6em, after-item-skip=5pt, before-skip=3pt, after-skip=4pt}

% ---------- capçalera i peu ----------
\newcommand{\llicencia}{\href{https://creativecommons.org/licenses/by-nc-sa/4.0/deed.ca}{\ccbyncsa}\ \ Llicència Creative Commons Reconeixement-NoComercial-CompartirIgual 4.0 Internacional (CC BY-NC-SA 4.0)}
\pagestyle{fancy}
\fancyhf{}
\fancyhead[L]{\raisebox{-0.35\height}{\includegraphics[height=0.95cm]{logo-gva.png}}}
\fancyhead[C]{\small\color{muted}Solucions de la fitxa · Unitat 6: Sistemes d'equacions · 3r d'ESO}
\fancyhead[R]{\raisebox{-0.35\height}{\includegraphics[height=0.95cm]{logo-vermellar.png}}}
\fancyfoot[L]{\scriptsize\color{muted}\llicencia}
\fancyfoot[R]{\small\color{muted}\thepage}
\renewcommand{\headrulewidth}{0.4pt}
\renewcommand{\headrule}{\color{muted!40}\hrule height \headrulewidth}
\fancypagestyle{primera}{\fancyhf{}\renewcommand{\headrulewidth}{0pt}
  \fancyfoot[L]{\scriptsize\color{muted}\llicencia}
  \fancyfoot[R]{\small\color{muted}\thepage}}

% ---------- nivells i exercicis ----------
\newcommand{\estrelles}[1]{{\color{accent}\ifcase#1\or\faStar\faStar[regular]\faStar[regular]\or\faStar\faStar\faStar[regular]\or\faStar\faStar\faStar\fi}}
\newcounter{exer}
\newcommand{\ex}[1]{\par\needspace{4\baselineskip}\medskip\stepcounter{exer}\noindent{\color{navy}\bfseries\theexer.}\ \ {\small\estrelles{#1}}\ \ }

\newcommand{\seccio}[2]{\par\needspace{10\baselineskip}\bigskip
  \begin{tcolorbox}[colback=navy,colframe=navy,coltext=white,arc=2mm,boxrule=0pt,left=3mm,right=3mm,top=1.2mm,bottom=1.2mm]
  {\large\bfseries #1\quad #2}
  \end{tcolorbox}\vspace{-2pt}}

\newcolumntype{C}[1]{>{\centering\arraybackslash}p{#1}}
\usepackage{pifont}

% quadre de recordatori (només a la fitxa d'exercicis)
\newcommand{\recorda}[1]{\par\smallskip\begin{tcolorbox}[colback=accent!7,colframe=accent!55,arc=2mm,boxrule=0.5pt,left=2.5mm,right=2.5mm,top=1mm,bottom=1mm]\small{\color{accent}\faLightbulb[regular]}\ \textbf{Recorda:}\ #1\end{tcolorbox}}
\newcommand{\res}[1]{\boxed{#1}}
\newcommand{\rest}[1]{\fbox{#1}}

\def\fmneed{14}
% =====================================================================
% Capçalera FentMates de les fitxes: un color per a cada curs.
% Els colors són els de la pàgina inicial (LEVELS a assets/js/materials.js);
% el logo ve de _eines/logo/pdf/ (python3 _eines/logo/genera.py).
%
% Ús, al final del preàmbul de la unitat:
%   % =====================================================================
% Capçalera FentMates de les fitxes: un color per a cada curs.
% Els colors són els de la pàgina inicial (LEVELS a assets/js/materials.js);
% el logo ve de _eines/logo/pdf/ (python3 _eines/logo/genera.py).
%
% Ús, al final del preàmbul de la unitat:
%   % =====================================================================
% Capçalera FentMates de les fitxes: un color per a cada curs.
% Els colors són els de la pàgina inicial (LEVELS a assets/js/materials.js);
% el logo ve de _eines/logo/pdf/ (python3 _eines/logo/genera.py).
%
% Ús, al final del preàmbul de la unitat:
%   \input{../comu/fentmates.tex}
%   \fmcurs{1eso}
% i, al començament del document:
%   \capcalera{Matemàtiques}{1r ESO}{Fitxa d'exercicis · Unitat 2}{Divisibilitat}
% Opcional abans de \input: \def\fmneed{14} (espai mínim abans de cada secció).
% =====================================================================
\makeatletter
\def\fm@logodir{../../logo/pdf}
\providecommand{\fmneed}{10}
\def\fm@cursnom{}

% clau del curs, color i color del text damunt del color
\newcommand{\fm@curs}[3]{\@namedef{fm@color@#1}{#2}\@namedef{fm@text@#1}{#3}}
\fm@curs{1eso}{E6B84A}{navy}
\fm@curs{2eso}{3B82F6}{white}
\fm@curs{3eso}{FF7A50}{white}
\fm@curs{4eso-a}{A78BFA}{white}
\fm@curs{4eso-b}{6D28D9}{white}
\fm@curs{1batx-mi}{9B1C1C}{white}
\fm@curs{1batx-ccss}{EF6B6B}{white}
\fm@curs{2batx-mii}{1E6B47}{white}
\fm@curs{2batx-ccss}{2BB3A3}{white}

\newcommand{\fmcurs}[1]{%
  \@ifundefined{fm@color@#1}{\PackageError{fentmates}{Curs desconegut: #1}{}}{}%
  \edef\fm@tmp{\noexpand\definecolor{curs}{HTML}{\@nameuse{fm@color@#1}}%
               \noexpand\colorlet{curstext}{\@nameuse{fm@text@#1}}}\fm@tmp
  \colorlet{cursfosc}{curs!72!black}%
  \def\fm@logo{\fm@logodir/vertical-#1.pdf}%
  \def\fm@marca{\fm@logodir/marca-#1.pdf}}

% ---------- capçalera de la primera pàgina ----------
% #1 matèria (petit), #2 curs, #3 tipus de document i unitat, #4 títol
\newcommand{\capcalera}[4]{%
  \gdef\fm@cursnom{#2}%
  \par\noindent
  \begin{tcolorbox}[colback=curs,colframe=curs,boxrule=0pt,arc=2.2mm,width=3.5cm,height=2.15cm,
      valign=center,left=3mm,right=2mm,top=1mm,bottom=1mm,nobeforeafter,box align=center]
    \color{curstext}{\footnotesize #1}\par\vspace{1pt}{\fontsize{19}{21}\selectfont\bfseries #2\par}%
  \end{tcolorbox}\hspace{5mm}%
  \begin{minipage}[c]{\dimexpr\linewidth-3.5cm-5mm-2.9cm\relax}
    {\color{cursfosc}\small\bfseries\MakeUppercase{#3}}\par\vspace{2pt}
    {\color{navy}\fontsize{20}{23}\selectfont\bfseries #4\par}
  \end{minipage}\hfill
  \begin{minipage}[c]{2.6cm}\raggedleft\includegraphics[width=2.6cm]{\fm@logo}\end{minipage}\par
  \vspace{2pt}}

% ---------- barres de secció del color del curs ----------
\renewcommand{\seccio}[2]{\par\needspace{\fmneed\baselineskip}\bigskip
  \begin{tcolorbox}[colback=curs,colframe=curs,coltext=curstext,arc=2mm,boxrule=0pt,left=3mm,right=3mm,top=1.2mm,bottom=1.2mm]
  {\large\bfseries #1\quad #2}
  \end{tcolorbox}\vspace{-2pt}}

% ---------- capçalera i peu de les pàgines ----------
% primera pàgina: logos institucionals al peu
\fancypagestyle{primera}{\fancyhf{}\renewcommand{\headrulewidth}{0pt}
  \fancyfoot[L]{\raisebox{-0.32\height}{\includegraphics[height=0.78cm]{logo-gva.png}}\hspace{5mm}%
    \raisebox{-0.32\height}{\includegraphics[height=0.78cm]{logo-vermellar.png}}}
  \fancyfoot[R]{\scriptsize\color{muted}\href{https://creativecommons.org/licenses/by-nc-sa/4.0/deed.ca}{\ccbyncsa}\ \ CC BY-NC-SA 4.0 · fentmates.org\qquad\small\thepage}}
% la resta: el símbol a l'esquerra i el curs a la dreta, amb el títol al centre com abans
\fancyhead[L]{\raisebox{-0.3\height}{\includegraphics[height=0.62cm]{\fm@marca}}}
\fancyhead[R]{\small\bfseries\color{cursfosc}\fm@cursnom}
\renewcommand{\headrule}{\color{curs!45}\hrule height \headrulewidth}
\makeatother

%   \fmcurs{1eso}
% i, al començament del document:
%   \capcalera{Matemàtiques}{1r ESO}{Fitxa d'exercicis · Unitat 2}{Divisibilitat}
% Opcional abans de \input: \def\fmneed{14} (espai mínim abans de cada secció).
% =====================================================================
\makeatletter
\def\fm@logodir{../../logo/pdf}
\providecommand{\fmneed}{10}
\def\fm@cursnom{}

% clau del curs, color i color del text damunt del color
\newcommand{\fm@curs}[3]{\@namedef{fm@color@#1}{#2}\@namedef{fm@text@#1}{#3}}
\fm@curs{1eso}{E6B84A}{navy}
\fm@curs{2eso}{3B82F6}{white}
\fm@curs{3eso}{FF7A50}{white}
\fm@curs{4eso-a}{A78BFA}{white}
\fm@curs{4eso-b}{6D28D9}{white}
\fm@curs{1batx-mi}{9B1C1C}{white}
\fm@curs{1batx-ccss}{EF6B6B}{white}
\fm@curs{2batx-mii}{1E6B47}{white}
\fm@curs{2batx-ccss}{2BB3A3}{white}

\newcommand{\fmcurs}[1]{%
  \@ifundefined{fm@color@#1}{\PackageError{fentmates}{Curs desconegut: #1}{}}{}%
  \edef\fm@tmp{\noexpand\definecolor{curs}{HTML}{\@nameuse{fm@color@#1}}%
               \noexpand\colorlet{curstext}{\@nameuse{fm@text@#1}}}\fm@tmp
  \colorlet{cursfosc}{curs!72!black}%
  \def\fm@logo{\fm@logodir/vertical-#1.pdf}%
  \def\fm@marca{\fm@logodir/marca-#1.pdf}}

% ---------- capçalera de la primera pàgina ----------
% #1 matèria (petit), #2 curs, #3 tipus de document i unitat, #4 títol
\newcommand{\capcalera}[4]{%
  \gdef\fm@cursnom{#2}%
  \par\noindent
  \begin{tcolorbox}[colback=curs,colframe=curs,boxrule=0pt,arc=2.2mm,width=3.5cm,height=2.15cm,
      valign=center,left=3mm,right=2mm,top=1mm,bottom=1mm,nobeforeafter,box align=center]
    \color{curstext}{\footnotesize #1}\par\vspace{1pt}{\fontsize{19}{21}\selectfont\bfseries #2\par}%
  \end{tcolorbox}\hspace{5mm}%
  \begin{minipage}[c]{\dimexpr\linewidth-3.5cm-5mm-2.9cm\relax}
    {\color{cursfosc}\small\bfseries\MakeUppercase{#3}}\par\vspace{2pt}
    {\color{navy}\fontsize{20}{23}\selectfont\bfseries #4\par}
  \end{minipage}\hfill
  \begin{minipage}[c]{2.6cm}\raggedleft\includegraphics[width=2.6cm]{\fm@logo}\end{minipage}\par
  \vspace{2pt}}

% ---------- barres de secció del color del curs ----------
\renewcommand{\seccio}[2]{\par\needspace{\fmneed\baselineskip}\bigskip
  \begin{tcolorbox}[colback=curs,colframe=curs,coltext=curstext,arc=2mm,boxrule=0pt,left=3mm,right=3mm,top=1.2mm,bottom=1.2mm]
  {\large\bfseries #1\quad #2}
  \end{tcolorbox}\vspace{-2pt}}

% ---------- capçalera i peu de les pàgines ----------
% primera pàgina: logos institucionals al peu
\fancypagestyle{primera}{\fancyhf{}\renewcommand{\headrulewidth}{0pt}
  \fancyfoot[L]{\raisebox{-0.32\height}{\includegraphics[height=0.78cm]{logo-gva.png}}\hspace{5mm}%
    \raisebox{-0.32\height}{\includegraphics[height=0.78cm]{logo-vermellar.png}}}
  \fancyfoot[R]{\scriptsize\color{muted}\href{https://creativecommons.org/licenses/by-nc-sa/4.0/deed.ca}{\ccbyncsa}\ \ CC BY-NC-SA 4.0 · fentmates.org\qquad\small\thepage}}
% la resta: el símbol a l'esquerra i el curs a la dreta, amb el títol al centre com abans
\fancyhead[L]{\raisebox{-0.3\height}{\includegraphics[height=0.62cm]{\fm@marca}}}
\fancyhead[R]{\small\bfseries\color{cursfosc}\fm@cursnom}
\renewcommand{\headrule}{\color{curs!45}\hrule height \headrulewidth}
\makeatother

%   \fmcurs{1eso}
% i, al començament del document:
%   \capcalera{Matemàtiques}{1r ESO}{Fitxa d'exercicis · Unitat 2}{Divisibilitat}
% Opcional abans de \input: \def\fmneed{14} (espai mínim abans de cada secció).
% =====================================================================
\makeatletter
\def\fm@logodir{../../logo/pdf}
\providecommand{\fmneed}{10}
\def\fm@cursnom{}

% clau del curs, color i color del text damunt del color
\newcommand{\fm@curs}[3]{\@namedef{fm@color@#1}{#2}\@namedef{fm@text@#1}{#3}}
\fm@curs{1eso}{E6B84A}{navy}
\fm@curs{2eso}{3B82F6}{white}
\fm@curs{3eso}{FF7A50}{white}
\fm@curs{4eso-a}{A78BFA}{white}
\fm@curs{4eso-b}{6D28D9}{white}
\fm@curs{1batx-mi}{9B1C1C}{white}
\fm@curs{1batx-ccss}{EF6B6B}{white}
\fm@curs{2batx-mii}{1E6B47}{white}
\fm@curs{2batx-ccss}{2BB3A3}{white}

\newcommand{\fmcurs}[1]{%
  \@ifundefined{fm@color@#1}{\PackageError{fentmates}{Curs desconegut: #1}{}}{}%
  \edef\fm@tmp{\noexpand\definecolor{curs}{HTML}{\@nameuse{fm@color@#1}}%
               \noexpand\colorlet{curstext}{\@nameuse{fm@text@#1}}}\fm@tmp
  \colorlet{cursfosc}{curs!72!black}%
  \def\fm@logo{\fm@logodir/vertical-#1.pdf}%
  \def\fm@marca{\fm@logodir/marca-#1.pdf}}

% ---------- capçalera de la primera pàgina ----------
% #1 matèria (petit), #2 curs, #3 tipus de document i unitat, #4 títol
\newcommand{\capcalera}[4]{%
  \gdef\fm@cursnom{#2}%
  \par\noindent
  \begin{tcolorbox}[colback=curs,colframe=curs,boxrule=0pt,arc=2.2mm,width=3.5cm,height=2.15cm,
      valign=center,left=3mm,right=2mm,top=1mm,bottom=1mm,nobeforeafter,box align=center]
    \color{curstext}{\footnotesize #1}\par\vspace{1pt}{\fontsize{19}{21}\selectfont\bfseries #2\par}%
  \end{tcolorbox}\hspace{5mm}%
  \begin{minipage}[c]{\dimexpr\linewidth-3.5cm-5mm-2.9cm\relax}
    {\color{cursfosc}\small\bfseries\MakeUppercase{#3}}\par\vspace{2pt}
    {\color{navy}\fontsize{20}{23}\selectfont\bfseries #4\par}
  \end{minipage}\hfill
  \begin{minipage}[c]{2.6cm}\raggedleft\includegraphics[width=2.6cm]{\fm@logo}\end{minipage}\par
  \vspace{2pt}}

% ---------- barres de secció del color del curs ----------
\renewcommand{\seccio}[2]{\par\needspace{\fmneed\baselineskip}\bigskip
  \begin{tcolorbox}[colback=curs,colframe=curs,coltext=curstext,arc=2mm,boxrule=0pt,left=3mm,right=3mm,top=1.2mm,bottom=1.2mm]
  {\large\bfseries #1\quad #2}
  \end{tcolorbox}\vspace{-2pt}}

% ---------- capçalera i peu de les pàgines ----------
% primera pàgina: logos institucionals al peu
\fancypagestyle{primera}{\fancyhf{}\renewcommand{\headrulewidth}{0pt}
  \fancyfoot[L]{\raisebox{-0.32\height}{\includegraphics[height=0.78cm]{logo-gva.png}}\hspace{5mm}%
    \raisebox{-0.32\height}{\includegraphics[height=0.78cm]{logo-vermellar.png}}}
  \fancyfoot[R]{\scriptsize\color{muted}\href{https://creativecommons.org/licenses/by-nc-sa/4.0/deed.ca}{\ccbyncsa}\ \ CC BY-NC-SA 4.0 · fentmates.org\qquad\small\thepage}}
% la resta: el símbol a l'esquerra i el curs a la dreta, amb el títol al centre com abans
\fancyhead[L]{\raisebox{-0.3\height}{\includegraphics[height=0.62cm]{\fm@marca}}}
\fancyhead[R]{\small\bfseries\color{cursfosc}\fm@cursnom}
\renewcommand{\headrule}{\color{curs!45}\hrule height \headrulewidth}
\makeatother

\fmcurs{3eso}

\begin{document}
\thispagestyle{primera}

\capcalera{Matemàtiques}{3r ESO}{Solucions de la fitxa · Unitat 6}{Sistemes d'equacions}
\vspace{8pt}

\begin{tcolorbox}[colback=soft,colframe=soft,arc=2mm,left=3mm,right=3mm,top=1.5mm,bottom=1.5mm]
\small Intenta resoldre cada exercici abans de mirar-ne la solució. Si no t'ix el mateix resultat, repassa els passos o consulta els apunts interactius de la unitat: hi trobaràs calculadores que resolen pas a pas qualsevol sistema pels tres mètodes i el dibuixen.
\end{tcolorbox}

% =====================================================
\seccio{1}{Equacions lineals amb dues incògnites}

\ex{1} \begin{tasks}(3)
  \task \rest{sí}
  \task no ($x^2$)
  \task no ($xy$)
  \task \rest{sí}
  \task no (una incògnita)
  \task \rest{sí}
\end{tasks}

\ex{1} \begin{tasks}(4)
  \task $6+6=12$: \rest{sí}
  \task $0+12=12$: \rest{sí}
  \task $12+3=15$: \rest{no}
  \task $-6+18=12$: \rest{sí}
\end{tasks}

\ex{1} \begin{tasks}(2)
  \task $y=5-x$: \ $y=7,\ 6,\ 5,\ 4,\ 3$
  \task $y=2x-1$: \ $y=-5,\ -3,\ -1,\ 1,\ 3$
\end{tasks}

\ex{2} \begin{tasks}(2)
  \task $x=6+3y$: \ p. ex.\ \rest{$(6,0),\ (9,1),\ (3,-1)$}
  \task $y=5-2x$: \ p. ex.\ \rest{$(0,5),\ (1,3),\ (2,1)$}
\end{tasks}

\ex{2} \begin{tasks}(1)
  \task $-4+3k=8\Rightarrow\res{k=4}$
  \task $10-2m=4\Rightarrow\res{m=3}$
\end{tasks}

\ex{2} \begin{tasks}(2)
  \task $(0,2),\ (2,1),\ (4,0)$\\[2pt]\begin{tikzpicture}[scale=0.34,baseline=(current bounding box.center)]
\draw[muted!25,very thin] (-5,-5) grid (5,5);
\draw[-latex,muted] (-5.4,0)--(5.6,0) node[right,font=\scriptsize]{$x$};
\draw[-latex,muted] (0,-5.4)--(0,5.6) node[above,font=\scriptsize]{$y$};
\node[font=\tiny,muted,below] at (-4,0) {-4}; \node[font=\tiny,muted,left] at (0,-4) {-4};
\node[font=\tiny,muted,below] at (-2,0) {-2}; \node[font=\tiny,muted,left] at (0,-2) {-2};
\node[font=\tiny,muted,below] at (2,0) {2}; \node[font=\tiny,muted,left] at (0,2) {2};
\node[font=\tiny,muted,below] at (4,0) {4}; \node[font=\tiny,muted,left] at (0,4) {4};
\begin{scope}\clip (-5,-5) rectangle (5,5);
\draw[accent,line width=1.1pt] (-6.000,5.000)--(6.000,-1.000);
\end{scope}
\end{tikzpicture}
  \task $(0,-3),\ (1,0),\ (2,3)$\\[2pt]\begin{tikzpicture}[scale=0.34,baseline=(current bounding box.center)]
\draw[muted!25,very thin] (-5,-5) grid (5,5);
\draw[-latex,muted] (-5.4,0)--(5.6,0) node[right,font=\scriptsize]{$x$};
\draw[-latex,muted] (0,-5.4)--(0,5.6) node[above,font=\scriptsize]{$y$};
\node[font=\tiny,muted,below] at (-4,0) {-4}; \node[font=\tiny,muted,left] at (0,-4) {-4};
\node[font=\tiny,muted,below] at (-2,0) {-2}; \node[font=\tiny,muted,left] at (0,-2) {-2};
\node[font=\tiny,muted,below] at (2,0) {2}; \node[font=\tiny,muted,left] at (0,2) {2};
\node[font=\tiny,muted,below] at (4,0) {4}; \node[font=\tiny,muted,left] at (0,4) {4};
\begin{scope}\clip (-5,-5) rectangle (5,5);
\draw[accent,line width=1.1pt] (-6.000,-21.000)--(6.000,15.000);
\end{scope}
\end{tikzpicture}
\end{tasks}

\ex{3} \begin{tasks}(1)
  \task $y=x+3$: \ $3x+2x+6=21\Rightarrow x=3$: \ \rest{$(3,\,6)$}
  \task Si $y=0$, $x=6$; si $x=0$, $y=-3$: \ \rest{$(6,\,0)$ i $(0,\,-3)$}
\end{tasks}


% =====================================================
\seccio{2}{Sistemes d'equacions i solucions}

\ex{1} \begin{tasks}(2)
  \task $1=1$ i $5=5$: \rest{sí}
  \task $5=5$, però $2+2=4\neq3$: \rest{no}
\end{tasks}

\ex{1} $(1,2)$: $2-2=0$ \ding{55}; \ $(2,1)$: $3$ i $3$ \checkmark; \ $(3,0)$: $6\neq3$ \ding{55}. \ La solució és \rest{$(2,\,1)$}.

\ex{2} Per exemple, \rest{$x+y=3$ i $x-y=-7$} (hi ha infinites respostes).

\ex{3} $3a-4=5\Rightarrow\res{a=3}$; \ $12-2b=6\Rightarrow\res{b=3}$


% =====================================================
\seccio{3}{Mètode de substitució}

\ex{1} \begin{tasks}(1)
  \task $x=3y$\newline $(3y)+y=8\Rightarrow 3y+y=8\Rightarrow 4y=8\Rightarrow y=\tfrac{8}{4}=2$\newline $x=3\cdot 2=6$\newline $\res{x=6,\ \ y=2}$
  \task $y=2x+1$\newline $3x+(2x+1)=11\Rightarrow 3x+2x+1=11\Rightarrow 5x=10\Rightarrow x=\tfrac{10}{5}=2$\newline $y=2\cdot 2+1=5$\newline $\res{x=2,\ \ y=5}$
\end{tasks}

\ex{1} \begin{tasks}(1)
  \task $x=7-y$\newline $2(7-y)-y=2\Rightarrow 14-2y-y=2\Rightarrow -3y=-12\Rightarrow y=\tfrac{-12}{-3}=4$\newline $x=7-4=3$\newline $\res{x=3,\ \ y=4}$
  \task $x=y+1$\newline $3(y+1)+2y=13\Rightarrow 3y+3+2y=13\Rightarrow 5y=10\Rightarrow y=\tfrac{10}{5}=2$\newline $x=2+1=3$\newline $\res{x=3,\ \ y=2}$
\end{tasks}

\ex{2} \begin{tasks}(1)
  \task $y=1-3x$\newline $5x-2(1-3x)=20\Rightarrow 5x-2+6x=20\Rightarrow 11x=22\Rightarrow x=\tfrac{22}{11}=2$\newline $y=1-3\cdot 2=-5$\newline $\res{x=2,\ \ y=-5}$
  \task $-y=-4-2x\Rightarrow y=2x+4$\newline $4x+3(2x+4)=2\Rightarrow 4x+6x+12=2\Rightarrow 10x=-10\Rightarrow x=\tfrac{-10}{10}=-1$\newline $y=2\cdot (-1)+4=2$\newline $\res{x=-1,\ \ y=2}$
\end{tasks}

\ex{2} \begin{tasks}(1)
  \task $x=-5-4y$\newline $3(-5-4y)-2y=13\Rightarrow -15-12y-2y=13\Rightarrow -14y=28\Rightarrow y=\tfrac{28}{-14}=-2$\newline $x=-5-4\cdot (-2)=3$\newline $\res{x=3,\ \ y=-2}$
  \task $-y=13-5x\Rightarrow y=5x-13$\newline $2x+3(5x-13)=12\Rightarrow 2x+15x-39=12\Rightarrow 17x=51\Rightarrow x=\tfrac{51}{17}=3$\newline $y=5\cdot 3-13=2$\newline $\res{x=3,\ \ y=2}$
\end{tasks}

\ex{3} \begin{tasks}(1)
  \task $2x=12-3y\Rightarrow x=\tfrac{12-3y}{2}$\newline $5\cdot \tfrac{12-3y}{2}-4y=7\Rightarrow 5(12-3y)-8y=14\Rightarrow 60-15y-8y=14\Rightarrow -23y=-46\Rightarrow y=\tfrac{-46}{-23}=2$\newline $x=\tfrac{12-3\cdot 2}{2}=3$\newline $\res{x=3,\ \ y=2}$
  \task $-2y=-1-3x\Rightarrow y=\tfrac{3x+1}{2}$\newline $4x+5\cdot \tfrac{3x+1}{2}=14\Rightarrow 8x+5(3x+1)=28\Rightarrow 8x+15x+5=28\Rightarrow 23x=23\Rightarrow x=\tfrac{23}{23}=1$\newline $y=\tfrac{3\cdot 1+1}{2}=2$\newline $\res{x=1,\ \ y=2}$
\end{tasks}

\ex{3} $-3(3-x)=-9+3x$: el $-3$ multiplica els dos termes. \ $y=3-x$\newline $2x-3(3-x)=1\Rightarrow 2x-9+3x=1\Rightarrow 5x=10\Rightarrow x=\tfrac{10}{5}=2$\newline $y=3-2=1$\newline $\res{x=2,\ \ y=1}$


% =====================================================
\seccio{4}{Mètode d'igualació}

\ex{1} \begin{tasks}(1)
  \task $y=x+5,\quad y=3x-1$\newline $x+5=3x-1\Rightarrow -2x=-6\Rightarrow x=\tfrac{-6}{-2}=3$\newline $y=3+5=8$\newline $\res{x=3,\ \ y=8}$
  \task $x=2y-3,\quad x=y+1$\newline $2y-3=y+1\Rightarrow y=4$\newline $x=2\cdot 4-3=5$\newline $\res{x=5,\ \ y=4}$
\end{tasks}

\ex{1} \begin{tasks}(1)
  \task $x=6-y,\quad x=y+2$\newline $6-y=y+2\Rightarrow -2y=-4\Rightarrow y=\tfrac{-4}{-2}=2$\newline $x=6-2=4$\newline $\res{x=4,\ \ y=2}$
  \task $y=2x+1,\quad y=7-x$\newline $2x+1=7-x\Rightarrow 3x=6\Rightarrow x=\tfrac{6}{3}=2$\newline $y=2\cdot 2+1=5$\newline $\res{x=2,\ \ y=5}$
\end{tasks}

\ex{2} \begin{tasks}(1)
  \task $y=8-2x,\quad y=11-3x$\newline $8-2x=11-3x\Rightarrow x=3$\newline $y=8-2\cdot 3=2$\newline $\res{x=3,\ \ y=2}$
  \task $x=10-3y,\quad x=2y$\newline $10-3y=2y\Rightarrow -5y=-10\Rightarrow y=\tfrac{-10}{-5}=2$\newline $x=10-3\cdot 2=4$\newline $\res{x=4,\ \ y=2}$
\end{tasks}

\ex{2} \begin{tasks}(1)
  \task $-y=7-3x\Rightarrow y=3x-7,\quad 2y=-x\Rightarrow y=\tfrac{-x}{2}$\newline $3x-7=\tfrac{-x}{2}\Rightarrow 2(3x-7)=-x\Rightarrow 6x-14=-x\Rightarrow 7x=14\Rightarrow x=\tfrac{14}{7}=2$\newline $y=3\cdot 2-7=-1$\newline $\res{x=2,\ \ y=-1}$
  \task $y=-2-4x,\quad -y=-4-2x\Rightarrow y=2x+4$\newline $-2-4x=2x+4\Rightarrow -6x=6\Rightarrow x=\tfrac{6}{-6}=-1$\newline $y=-2-4\cdot (-1)=2$\newline $\res{x=-1,\ \ y=2}$
\end{tasks}

\ex{3} \begin{tasks}(1)
  \task $2x=7-3y\Rightarrow x=\tfrac{7-3y}{2},\quad 3x=2y+4\Rightarrow x=\tfrac{2y+4}{3}$\newline $\tfrac{7-3y}{2}=\tfrac{2y+4}{3}\Rightarrow 3(7-3y)=2(2y+4)\Rightarrow 21-9y=4y+8\Rightarrow -13y=-13\Rightarrow y=\tfrac{-13}{-13}=1$\newline $x=\tfrac{7-3\cdot 1}{2}=2$\newline $\res{x=2,\ \ y=1}$
  \task $-2y=1-5x\Rightarrow y=\tfrac{5x-1}{2},\quad 4y=11-3x\Rightarrow y=\tfrac{11-3x}{4}$\newline $\tfrac{5x-1}{2}=\tfrac{11-3x}{4}\Rightarrow 2(5x-1)=11-3x\Rightarrow 10x-2=11-3x\Rightarrow 13x=13\Rightarrow x=\tfrac{13}{13}=1$\newline $y=\tfrac{5\cdot 1-1}{2}=2$\newline $\res{x=1,\ \ y=2}$
\end{tasks}

\ex{3} Igualació: $\tfrac{x+1}{2}=\tfrac{2x-1}{3}\Rightarrow 3(x+1)=2(2x-1)\Rightarrow 3x+3=4x-2\Rightarrow x=5$; \ $y=\tfrac{5+1}{2}=3$. \ $\res{x=5,\ \ y=3}$


% =====================================================
\seccio{5}{Mètode de reducció}

\ex{1} \begin{tasks}(1)
  \task $\begin{array}{lrrcr}&x&+y&=&9\\&x&-y&=&3\\\hline&2x&&=&12\end{array}\Rightarrow 2x=12\Rightarrow x=\tfrac{12}{2}=6$\newline $6+y=9\Rightarrow y=3$\newline $\res{x=6,\ \ y=3}$
  \task $\begin{array}{lrrcr}&3x&+2y&=&13\\&x&-2y&=&-1\\\hline&4x&&=&12\end{array}\Rightarrow 4x=12\Rightarrow x=\tfrac{12}{4}=3$\newline $3\cdot 3+2y=13\Rightarrow 2y=4\Rightarrow y=\tfrac{4}{2}=2$\newline $\res{x=3,\ \ y=2}$
\end{tasks}

\ex{1} \begin{tasks}(1)
  \task $\begin{array}{lrrcr}&2x&+y&=&7\\\scriptstyle(\cdot(-1))&-2x&+3y&=&5\\\hline&&4y&=&12\end{array}\Rightarrow 4y=12\Rightarrow y=\tfrac{12}{4}=3$\newline $2x+3=7\Rightarrow 2x=4\Rightarrow x=\tfrac{4}{2}=2$\newline $\res{x=2,\ \ y=3}$
  \task $\begin{array}{lrrcr}&5x&+3y&=&11\\&2x&-3y&=&-4\\\hline&7x&&=&7\end{array}\Rightarrow 7x=7\Rightarrow x=\tfrac{7}{7}=1$\newline $5\cdot 1+3y=11\Rightarrow 3y=6\Rightarrow y=\tfrac{6}{3}=2$\newline $\res{x=1,\ \ y=2}$
\end{tasks}

\ex{2} \begin{tasks}(1)
  \task $\begin{array}{lrrcr}&x&+2y&=&4\\\scriptstyle(\cdot2)&6x&-2y&=&10\\\hline&7x&&=&14\end{array}\Rightarrow 7x=14\Rightarrow x=\tfrac{14}{7}=2$\newline $3\cdot 2-y=5\Rightarrow -y=-1\Rightarrow y=\tfrac{-1}{-1}=1$\newline $\res{x=2,\ \ y=1}$
  \task $\begin{array}{lrrcr}\scriptstyle(\cdot(-1))&-4x&-3y&=&-10\\\scriptstyle(\cdot2)&4x&+10y&=&24\\\hline&&7y&=&14\end{array}\Rightarrow 7y=14\Rightarrow y=\tfrac{14}{7}=2$\newline $2x+5\cdot 2=12\Rightarrow 2x=2\Rightarrow x=\tfrac{2}{2}=1$\newline $\res{x=1,\ \ y=2}$
\end{tasks}

\ex{2} \begin{tasks}(1)
  \task $\begin{array}{lrrcr}&3x&+4y&=&18\\\scriptstyle(\cdot2)&10x&-4y&=&8\\\hline&13x&&=&26\end{array}\Rightarrow 13x=26\Rightarrow x=\tfrac{26}{13}=2$\newline $5\cdot 2-2y=4\Rightarrow -2y=-6\Rightarrow y=\tfrac{-6}{-2}=3$\newline $\res{x=2,\ \ y=3}$
  \task $\begin{array}{lrrcr}\scriptstyle(\cdot3)&6x&-15y&=&-3\\\scriptstyle(\cdot(-2))&-6x&-4y&=&-16\\\hline&&-19y&=&-19\end{array}\Rightarrow -19y=-19\Rightarrow y=\tfrac{-19}{-19}=1$\newline $2x-5\cdot 1=-1\Rightarrow 2x=4\Rightarrow x=\tfrac{4}{2}=2$\newline $\res{x=2,\ \ y=1}$
\end{tasks}

\ex{3} \begin{tasks}(1)
  \task $\begin{array}{lrrcr}\scriptstyle(\cdot2)&8x&+6y&=&-2\\\scriptstyle(\cdot(-3))&-15x&-6y&=&-12\\\hline&-7x&&=&-14\end{array}\Rightarrow -7x=-14\Rightarrow x=\tfrac{-14}{-7}=2$\newline $5\cdot 2+2y=4\Rightarrow 2y=-6\Rightarrow y=\tfrac{-6}{2}=-3$\newline $\res{x=2,\ \ y=-3}$
  \task $\begin{array}{lrrcr}\scriptstyle(\cdot2)&12x&-10y&=&14\\\scriptstyle(\cdot(-3))&-12x&+9y&=&-15\\\hline&&-y&=&-1\end{array}\Rightarrow -y=-1\Rightarrow y=\tfrac{-1}{-1}=1$\newline $4x-3\cdot 1=5\Rightarrow 4x=8\Rightarrow x=\tfrac{8}{4}=2$\newline $\res{x=2,\ \ y=1}$
\end{tasks}

\ex{3} $(\cdot2),\ (\cdot3)$: $13x=11\Rightarrow\res{x=\tfrac{11}{13}}$; \ $(\cdot3),\ (\cdot(-2))$: $13y=10\Rightarrow\res{y=\tfrac{10}{13}}$. \ Així evitem substituir una fracció.


% =====================================================
\seccio{6}{Mètode gràfic}

\ex{1} \begin{tasks}(2)
  \task Les rectes es tallen en $(1,\,3)$: \rest{$x=1,\ y=3$}
  \task Les rectes es tallen en $(1,\,-1)$: \rest{$x=1,\ y=-1$}
\end{tasks}

\ex{1} \begin{tasks}(2)
  \task \begin{tikzpicture}[scale=0.34,baseline=(current bounding box.center)]
\draw[muted!25,very thin] (-5,-5) grid (5,5);
\draw[-latex,muted] (-5.4,0)--(5.6,0) node[right,font=\scriptsize]{$x$};
\draw[-latex,muted] (0,-5.4)--(0,5.6) node[above,font=\scriptsize]{$y$};
\node[font=\tiny,muted,below] at (-4,0) {-4}; \node[font=\tiny,muted,left] at (0,-4) {-4};
\node[font=\tiny,muted,below] at (-2,0) {-2}; \node[font=\tiny,muted,left] at (0,-2) {-2};
\node[font=\tiny,muted,below] at (2,0) {2}; \node[font=\tiny,muted,left] at (0,2) {2};
\node[font=\tiny,muted,below] at (4,0) {4}; \node[font=\tiny,muted,left] at (0,4) {4};
\begin{scope}\clip (-5,-5) rectangle (5,5);
\draw[accent,line width=1.1pt] (-6.000,12.000)--(6.000,0.000);
\draw[teal,line width=1.1pt] (-6.000,-8.000)--(6.000,4.000);
\end{scope}
\fill[navy] (4,2) circle (0.22);
\node[font=\scriptsize\bfseries,navy,fill=white,inner sep=1pt,above right] at (4,2) {$(4,2)$};
\end{tikzpicture}\\[2pt] \rest{$x=4,\ y=2$}
  \task \begin{tikzpicture}[scale=0.34,baseline=(current bounding box.center)]
\draw[muted!25,very thin] (-5,-5) grid (5,5);
\draw[-latex,muted] (-5.4,0)--(5.6,0) node[right,font=\scriptsize]{$x$};
\draw[-latex,muted] (0,-5.4)--(0,5.6) node[above,font=\scriptsize]{$y$};
\node[font=\tiny,muted,below] at (-4,0) {-4}; \node[font=\tiny,muted,left] at (0,-4) {-4};
\node[font=\tiny,muted,below] at (-2,0) {-2}; \node[font=\tiny,muted,left] at (0,-2) {-2};
\node[font=\tiny,muted,below] at (2,0) {2}; \node[font=\tiny,muted,left] at (0,2) {2};
\node[font=\tiny,muted,below] at (4,0) {4}; \node[font=\tiny,muted,left] at (0,4) {4};
\begin{scope}\clip (-5,-5) rectangle (5,5);
\draw[accent,line width=1.1pt] (-6.000,-5.000)--(6.000,7.000);
\draw[teal,line width=1.1pt] (-6.000,11.000)--(6.000,-1.000);
\end{scope}
\fill[navy] (2,3) circle (0.22);
\node[font=\scriptsize\bfseries,navy,fill=white,inner sep=1pt,above right] at (2,3) {$(2,3)$};
\end{tikzpicture}\\[2pt] \rest{$x=2,\ y=3$}
\end{tasks}

\ex{2} \begin{tasks}(2)
  \task \begin{tikzpicture}[scale=0.34,baseline=(current bounding box.center)]
\draw[muted!25,very thin] (-5,-5) grid (5,5);
\draw[-latex,muted] (-5.4,0)--(5.6,0) node[right,font=\scriptsize]{$x$};
\draw[-latex,muted] (0,-5.4)--(0,5.6) node[above,font=\scriptsize]{$y$};
\node[font=\tiny,muted,below] at (-4,0) {-4}; \node[font=\tiny,muted,left] at (0,-4) {-4};
\node[font=\tiny,muted,below] at (-2,0) {-2}; \node[font=\tiny,muted,left] at (0,-2) {-2};
\node[font=\tiny,muted,below] at (2,0) {2}; \node[font=\tiny,muted,left] at (0,2) {2};
\node[font=\tiny,muted,below] at (4,0) {4}; \node[font=\tiny,muted,left] at (0,4) {4};
\begin{scope}\clip (-5,-5) rectangle (5,5);
\draw[accent,line width=1.1pt] (-6.000,16.000)--(6.000,-8.000);
\draw[teal,line width=1.1pt] (-6.000,-5.000)--(6.000,7.000);
\end{scope}
\fill[navy] (1,2) circle (0.22);
\node[font=\scriptsize\bfseries,navy,fill=white,inner sep=1pt,above right] at (1,2) {$(1,2)$};
\end{tikzpicture}\\[2pt] \rest{$x=1,\ y=2$}
  \task \begin{tikzpicture}[scale=0.34,baseline=(current bounding box.center)]
\draw[muted!25,very thin] (-5,-5) grid (5,5);
\draw[-latex,muted] (-5.4,0)--(5.6,0) node[right,font=\scriptsize]{$x$};
\draw[-latex,muted] (0,-5.4)--(0,5.6) node[above,font=\scriptsize]{$y$};
\node[font=\tiny,muted,below] at (-4,0) {-4}; \node[font=\tiny,muted,left] at (0,-4) {-4};
\node[font=\tiny,muted,below] at (-2,0) {-2}; \node[font=\tiny,muted,left] at (0,-2) {-2};
\node[font=\tiny,muted,below] at (2,0) {2}; \node[font=\tiny,muted,left] at (0,2) {2};
\node[font=\tiny,muted,below] at (4,0) {4}; \node[font=\tiny,muted,left] at (0,4) {4};
\begin{scope}\clip (-5,-5) rectangle (5,5);
\draw[accent,line width=1.1pt] (-6.000,3.000)--(6.000,-3.000);
\draw[teal,line width=1.1pt] (-6.000,-25.000)--(6.000,11.000);
\end{scope}
\fill[navy] (2,-1) circle (0.22);
\node[font=\scriptsize\bfseries,navy,fill=white,inner sep=1pt,above right] at (2,-1) {$(2,-1)$};
\end{tikzpicture}\\[2pt] \rest{$x=2,\ y=-1$}
\end{tasks}

\ex{2} \begin{tasks}(2)
  \task \begin{tikzpicture}[scale=0.34,baseline=(current bounding box.center)]
\draw[muted!25,very thin] (-5,-5) grid (5,5);
\draw[-latex,muted] (-5.4,0)--(5.6,0) node[right,font=\scriptsize]{$x$};
\draw[-latex,muted] (0,-5.4)--(0,5.6) node[above,font=\scriptsize]{$y$};
\node[font=\tiny,muted,below] at (-4,0) {-4}; \node[font=\tiny,muted,left] at (0,-4) {-4};
\node[font=\tiny,muted,below] at (-2,0) {-2}; \node[font=\tiny,muted,left] at (0,-2) {-2};
\node[font=\tiny,muted,below] at (2,0) {2}; \node[font=\tiny,muted,left] at (0,2) {2};
\node[font=\tiny,muted,below] at (4,0) {4}; \node[font=\tiny,muted,left] at (0,4) {4};
\begin{scope}\clip (-5,-5) rectangle (5,5);
\draw[accent,line width=1.1pt] (-6.000,-16.000)--(6.000,8.000);
\draw[teal,line width=1.1pt] (-6.000,5.000)--(6.000,-7.000);
\end{scope}
\fill[navy] (1,-2) circle (0.22);
\node[font=\scriptsize\bfseries,navy,fill=white,inner sep=1pt,above right] at (1,-2) {$(1,-2)$};
\end{tikzpicture}\\[2pt] \rest{$x=1,\ y=-2$}
  \task \begin{tikzpicture}[scale=0.34,baseline=(current bounding box.center)]
\draw[muted!25,very thin] (-5,-5) grid (5,5);
\draw[-latex,muted] (-5.4,0)--(5.6,0) node[right,font=\scriptsize]{$x$};
\draw[-latex,muted] (0,-5.4)--(0,5.6) node[above,font=\scriptsize]{$y$};
\node[font=\tiny,muted,below] at (-4,0) {-4}; \node[font=\tiny,muted,left] at (0,-4) {-4};
\node[font=\tiny,muted,below] at (-2,0) {-2}; \node[font=\tiny,muted,left] at (0,-2) {-2};
\node[font=\tiny,muted,below] at (2,0) {2}; \node[font=\tiny,muted,left] at (0,2) {2};
\node[font=\tiny,muted,below] at (4,0) {4}; \node[font=\tiny,muted,left] at (0,4) {4};
\begin{scope}\clip (-5,-5) rectangle (5,5);
\draw[accent,line width=1.1pt] (-6.000,-3.000)--(6.000,1.000);
\draw[teal,line width=1.1pt] (-6.000,4.000)--(6.000,-20.000);
\end{scope}
\fill[navy] (-3,-2) circle (0.22);
\node[font=\scriptsize\bfseries,navy,fill=white,inner sep=1pt,above right] at (-3,-2) {$(-3,-2)$};
\end{tikzpicture}\\[2pt] \rest{$x=-3,\ y=-2$}
\end{tasks}

\ex{3} \begin{tasks}(2)
  \task \begin{tikzpicture}[scale=0.34,baseline=(current bounding box.center)]
\draw[muted!25,very thin] (-5,-5) grid (5,5);
\draw[-latex,muted] (-5.4,0)--(5.6,0) node[right,font=\scriptsize]{$x$};
\draw[-latex,muted] (0,-5.4)--(0,5.6) node[above,font=\scriptsize]{$y$};
\node[font=\tiny,muted,below] at (-4,0) {-4}; \node[font=\tiny,muted,left] at (0,-4) {-4};
\node[font=\tiny,muted,below] at (-2,0) {-2}; \node[font=\tiny,muted,left] at (0,-2) {-2};
\node[font=\tiny,muted,below] at (2,0) {2}; \node[font=\tiny,muted,left] at (0,2) {2};
\node[font=\tiny,muted,below] at (4,0) {4}; \node[font=\tiny,muted,left] at (0,4) {4};
\begin{scope}\clip (-5,-5) rectangle (5,5);
\draw[accent,line width=1.1pt] (-6.000,8.000)--(6.000,-4.000);
\draw[teal,line width=1.1pt] (-6.000,9.000)--(6.000,-3.000);
\end{scope}
\end{tikzpicture}\\[2pt] Rectes paral·leles: \rest{cap solució}
  \task \begin{tikzpicture}[scale=0.34,baseline=(current bounding box.center)]
\draw[muted!25,very thin] (-5,-5) grid (5,5);
\draw[-latex,muted] (-5.4,0)--(5.6,0) node[right,font=\scriptsize]{$x$};
\draw[-latex,muted] (0,-5.4)--(0,5.6) node[above,font=\scriptsize]{$y$};
\node[font=\tiny,muted,below] at (-4,0) {-4}; \node[font=\tiny,muted,left] at (0,-4) {-4};
\node[font=\tiny,muted,below] at (-2,0) {-2}; \node[font=\tiny,muted,left] at (0,-2) {-2};
\node[font=\tiny,muted,below] at (2,0) {2}; \node[font=\tiny,muted,left] at (0,2) {2};
\node[font=\tiny,muted,below] at (4,0) {4}; \node[font=\tiny,muted,left] at (0,4) {4};
\begin{scope}\clip (-5,-5) rectangle (5,5);
\draw[accent,line width=1.1pt] (-6.000,-4.000)--(6.000,2.000);
\draw[teal,dashed,line width=1.6pt] (-6.000,-4.000)--(6.000,2.000);
\end{scope}
\end{tikzpicture}\\[2pt] La mateixa recta: \rest{infinites solucions}
\end{tasks}


% =====================================================
\seccio{7}{Classificació dels sistemes}

\ex{1} \begin{tasks}(2)
  \task $\tfrac{1}{2}=\tfrac{3}{6}=\tfrac{2}{4}$: \rest{SCI}
  \task $\tfrac{2}{4}=\tfrac{-1}{-2}\neq\tfrac{1}{3}$: \rest{SI}
  \task $\tfrac{1}{2}\neq\tfrac{1}{-1}$: \rest{SCD}
  \task $\tfrac{3}{-1}=\tfrac{-6}{2}=\tfrac{9}{-3}$: \rest{SCI}
\end{tasks}

\ex{1} \begin{tasks}(1)
  \task \rest{una solució}, \rest{compatible determinat}
  \task \rest{cap solució}, \rest{incompatible}
  \task \rest{infinites solucions}, \rest{compatible indeterminat}
\end{tasks}

\ex{2} \begin{tasks}(1)
  \task $\begin{array}{lrrcr}\scriptstyle(\cdot2)&2x&+2y&=&6\\\scriptstyle(\cdot(-1))&-2x&-2y&=&-5\\\hline&&0&=&1\end{array}$\newline \rest{SI} (cap solució)
  \task $\begin{array}{lrrcr}&x&-2y&=&1\\\scriptstyle(\cdot(-1))&-x&+2y&=&-1\\\hline&&0&=&0\end{array}$\newline \rest{SCI} (infinites solucions). \ Per exemple, $(1,\,0),\ (3,\,1),\ (-1,\,-1)$
\end{tasks}

\ex{2} \begin{tasks}(1)
  \task $\begin{array}{lrrcr}&2x&+3y&=&5\\\scriptstyle(\cdot(-1))&-2x&-3y&=&-5\\\hline&&0&=&0\end{array}$\newline \rest{SCI} (infinites solucions). \ P. ex.\ $(1,\,1),\ (4,\,-1)$
  \task $\begin{array}{lrrcr}&x&-y&=&2\\&-x&+y&=&2\\\hline&&0&=&4\end{array}$\newline \rest{SI} (cap solució)
  \task $\begin{array}{lrrcr}&3x&+y&=&2\\&x&-y&=&2\\\hline&4x&&=&4\end{array}\Rightarrow 4x=4\Rightarrow x=\tfrac{4}{4}=1$\newline $3\cdot 1+y=2\Rightarrow y=-1$\newline $\res{x=1,\ \ y=-1}$ \ (\rest{SCD})
\end{tasks}

\ex{2} \begin{tasks}(3)
  \task P. ex.\ \rest{$x+y=3$}
  \task P. ex.\ \rest{$4x-2y=6$}
  \task P. ex.\ \rest{$2x-y=5$}
\end{tasks}

\ex{3} $\tfrac m3=\tfrac21\Rightarrow\res{m=6}$; \ i aleshores $\tfrac31=3\neq2$: \rest{incompatible}.\newline Si $m\neq6$: \rest{compatible determinat}.

\ex{3} $\tfrac24=\tfrac{a}{-6}=\tfrac{b}{10}=\tfrac12\Rightarrow\res{a=-3,\ \ b=5}$


% =====================================================
\seccio{8}{Sistemes amb parèntesis i denominadors}

\ex{1} \begin{tasks}(1)
  \task $2(x-1)+y=7\Rightarrow 2x+y=9$\newline $x-(y-3)=3\Rightarrow x-y=0$\newline $y=9-2x$\newline $x-(9-2x)=0\Rightarrow x-9+2x=0\Rightarrow 3x=9\Rightarrow x=\tfrac{9}{3}=3$\newline $y=9-2\cdot 3=3$\newline $\res{x=3,\ \ y=3}$
  \task $3(x+y)-2x=10\Rightarrow x+3y=10$\newline $2(x-y)+y=-1\Rightarrow 2x-y=-1$\newline $x=10-3y$\newline $2(10-3y)-y=-1\Rightarrow 20-6y-y=-1\Rightarrow -7y=-21\Rightarrow y=\tfrac{-21}{-7}=3$\newline $x=10-3\cdot 3=1$\newline $\res{x=1,\ \ y=3}$
\end{tasks}

\ex{1} \begin{tasks}(1)
  \task $\dfrac{x}{2}+y=4\Rightarrow \cdot2:\ x+2y=8\Rightarrow x+2y=8$\newline $x-\dfrac{y}{3}=1\Rightarrow \cdot3:\ 3x-y=3\Rightarrow 3x-y=3$\newline $x=8-2y$\newline $3(8-2y)-y=3\Rightarrow 24-6y-y=3\Rightarrow -7y=-21\Rightarrow y=\tfrac{-21}{-7}=3$\newline $x=8-2\cdot 3=2$\newline $\res{x=2,\ \ y=3}$
  \task $\dfrac{x}{3}+\dfrac{y}{2}=2\Rightarrow \cdot6:\ 2x+3y=12\Rightarrow 2x+3y=12$\newline $\dfrac{x}{2}-\dfrac{y}{4}=1\Rightarrow \cdot4:\ 2x-y=4\Rightarrow 2x-y=4$\newline $\begin{array}{lrrcr}&2x&+3y&=&12\\\scriptstyle(\cdot(-1))&-2x&+y&=&-4\\\hline&&4y&=&8\end{array}\Rightarrow 4y=8\Rightarrow y=\tfrac{8}{4}=2$\newline $2x+3\cdot 2=12\Rightarrow 2x=6\Rightarrow x=\tfrac{6}{2}=3$\newline $\res{x=3,\ \ y=2}$
\end{tasks}

\ex{2} \begin{tasks}(1)
  \task $\dfrac{x+1}{2}+\dfrac{y-1}{3}=3\Rightarrow \cdot6:\ 3(x+1)+2(y-1)=18\Rightarrow 3x+2y=17$\newline $x=y-1$\newline $3(y-1)+2y=17\Rightarrow 3y-3+2y=17\Rightarrow 5y=20\Rightarrow y=\tfrac{20}{5}=4$\newline $x=4-1=3$\newline $\res{x=3,\ \ y=4}$
  \task $\dfrac{x-2}{3}=\dfrac{y+1}{4}\Rightarrow \cdot12:\ 4(x-2)=3(y+1)\Rightarrow 4x-3y=11$\newline $-y=7-2x\Rightarrow y=2x-7$\newline $4x-3(2x-7)=11\Rightarrow 4x-6x+21=11\Rightarrow -2x=-10\Rightarrow x=\tfrac{-10}{-2}=5$\newline $y=2\cdot 5-7=3$\newline $\res{x=5,\ \ y=3}$
\end{tasks}

\ex{2} $0{,}5x+0{,}2y=2\Rightarrow 5x+2y=20$\newline $x=y-3$\newline $5(y-3)+2y=20\Rightarrow 5y-15+2y=20\Rightarrow 7y=35\Rightarrow y=\tfrac{35}{7}=5$\newline $x=5-3=2$\newline $\res{x=2,\ \ y=5}$

\ex{3} $\dfrac{2x-1}{3}-\dfrac{y+2}{2}=-1\Rightarrow \cdot6:\ 2(2x-1)-3(y+2)=-6\Rightarrow 4x-3y=2$\newline $3(x-y)-2(x+1)=-6\Rightarrow x-3y=-4$\newline $\begin{array}{lrrcr}&4x&-3y&=&2\\\scriptstyle(\cdot(-1))&-x&+3y&=&4\\\hline&3x&&=&6\end{array}\Rightarrow 3x=6\Rightarrow x=\tfrac{6}{3}=2$\newline $4\cdot 2-3y=2\Rightarrow -3y=-6\Rightarrow y=\tfrac{-6}{-3}=2$\newline $\res{x=2,\ \ y=2}$

\ex{3} \begin{tasks}(1)
  \task Igualació (la $y$ ja està aïllada): $y=3x-2,\quad y=x+4$\newline $3x-2=x+4\Rightarrow 2x=6\Rightarrow x=\tfrac{6}{2}=3$\newline $y=3\cdot 3-2=7$\newline $\res{x=3,\ \ y=7}$
  \task Substitució ($x$ té coeficient 1): $x=9-5y$\newline $2(9-5y)-3y=5\Rightarrow 18-10y-3y=5\Rightarrow -13y=-13\Rightarrow y=\tfrac{-13}{-13}=1$\newline $x=9-5\cdot 1=4$\newline $\res{x=4,\ \ y=1}$
  \task Reducció (els coeficients de la $y$ són oposats): $\begin{array}{lrrcr}\scriptstyle(\cdot(-1))&-4x&-3y&=&-6\\\scriptstyle(\cdot2)&4x&-2y&=&16\\\hline&&-5y&=&10\end{array}\Rightarrow -5y=10\Rightarrow y=\tfrac{10}{-5}=-2$\newline $2x-(-2)=8\Rightarrow 2x=6\Rightarrow x=\tfrac{6}{2}=3$\newline $\res{x=3,\ \ y=-2}$
\end{tasks}


% =====================================================
\seccio{9}{Problemes amb sistemes d'equacions}

\ex{1} $\begin{cases}x+y=56\\ x-y=12\end{cases}\Rightarrow 2x=68\Rightarrow x=34,\ y=22$: \ \rest{34 i 22}

\ex{1} $\begin{cases}x=3y\\ x+y=48\end{cases}\Rightarrow 4y=48\Rightarrow y=12,\ x=36$: \ \rest{36 i 12}

\ex{1} $\begin{cases}x+y=40\\ 2x+4y=112\end{cases}\Rightarrow 2y=32\Rightarrow y=16,\ x=24$: \ \rest{24 gallines i 16 conills}

\ex{1} $\begin{cases}3x+2y=8\\ x+3y=5\end{cases}\Rightarrow x=5-3y\Rightarrow 15-7y=8\Rightarrow y=1,\ x=2$: \ \rest{pomes 2 €/kg, peres 1 €/kg}

\ex{2} $\begin{cases}x+y=22\\ 2x+0{,}5y=26\end{cases}\Rightarrow 1{,}5x=15\Rightarrow x=10,\ y=12$: \ \rest{10 de 2 € i 12 de 50 cèntims}

\ex{2} $\begin{cases}x+y=85\\ 7x+4y=445\end{cases}\Rightarrow 3x=105\Rightarrow x=35,\ y=50$: \ \rest{35 d'adult i 50 de xiquet}

\ex{2} $\begin{cases}x=4y\\ x+5=3(y+5)\end{cases}\Rightarrow 4y+5=3y+15\Rightarrow y=10$: \ \rest{mare 40 anys, fill 10 anys}

\ex{2} $\begin{cases}2x+2y=46\\ x-y=7\end{cases}\Rightarrow x+y=23\Rightarrow 2x=30$: \ \rest{base 15 m, altura 8 m}

\ex{2} $\begin{cases}x+y=30\\ 2x-y=39\end{cases}\Rightarrow 3x=69\Rightarrow x=23,\ y=7$: \ \rest{23 encerts (i 7 errades)}

\ex{2} $\begin{cases}x+y=28\\ y+2=x\end{cases}\Rightarrow 2y+2=28\Rightarrow y=13$: \ \rest{15 xics i 13 xiques}

\ex{3} $\begin{cases}x+y=60\\ 9x+6y=420\end{cases}\Rightarrow 3x=60\Rightarrow x=20,\ y=40$: \ \rest{20 kg de 9 € i 40 kg de 6 €}

\ex{3} $\begin{cases}x+y=11\\ 10y+x=10x+y-27\end{cases}\Rightarrow\begin{cases}x+y=11\\ x-y=3\end{cases}\Rightarrow x=7,\ y=4$: \ \rest{74}

\ex{3} $\begin{cases}3x+3y=75\\ x-y=5\end{cases}\Rightarrow x+y=25\Rightarrow x=15,\ y=10$: \ \rest{15 km/h i 10 km/h}

\ex{3} $\begin{cases}3(x+1)=2(y+1)\\ 2(x-1)=y-1\end{cases}\Rightarrow\begin{cases}3x-2y=-1\\ 2x-y=1\end{cases}\Rightarrow x=3,\ y=5$: \ \rest{$\tfrac35$}

\ex{3} $\begin{cases}x+y=50\\ 2x+y=87\end{cases}\Rightarrow x=37,\ y=13$: \ \rest{37 dobles i 13 individuals}

\ex{3} $\begin{cases}2x+y=40\\ x=y+5\end{cases}\Rightarrow 3y+10=40\Rightarrow y=10$: \ \rest{costats iguals 15 cm, base 10 cm}


\end{document}
